% 图契约：环中的支路值与端口测量值不同；图中的端口、几何、箭头和方程对应本题。
% 版面：沿用前三模块的几何图+推导分区；目标印宽17cm，简图至中等密度。
% 布局：左侧真实结构，右侧关键关系；真实电路边不作为架构流程箭头处理。
\documentclass[tikz,border=7pt]{standalone}
\usepackage{fontspec,xeCJK,amsmath,amssymb}
\IfFontExistsTF{Times New Roman}{\setmainfont{Times New Roman}}{\setmainfont{TeX Gyre Termes}}
\IfFontExistsTF{Heiti SC}{\setCJKmainfont[AutoFakeBold=1.5]{Heiti SC}}{\setCJKmainfont[AutoFakeBold=1.5]{FandolHei-Regular}}
\usetikzlibrary{arrows.meta,bending,calc,patterns,positioning,decorations.pathmorphing}
\definecolor{ink}{HTML}{243745}
\definecolor{blue}{HTML}{275DAB}
\definecolor{green}{HTML}{177C63}
\definecolor{red}{HTML}{BD4935}
\definecolor{orange}{HTML}{AA6C16}
\definecolor{purple}{HTML}{715393}
\tikzset{every node/.style={font=\fontsize{9.3}{12.5}\selectfont,text=ink},
arrow/.style={-{Stealth[length=5pt 1.5,width'=0pt 3]},shorten >=2pt,shorten <=1pt,line width=1pt},
arrow short/.style={-{Stealth[length=3pt 1.0,width'=0pt 2.5]},shorten >=1pt,shorten <=0pt,line width=.8pt},
axis/.style={arrow,draw=ink!60,line width=.6pt},
vec/.style={arrow,draw=blue},
wire/.style={draw=ink!75,line width=.8pt},
guide/.style={dashed,draw=ink!45,line width=.5pt},
card/.style={rounded corners=2pt,draw=blue!65,fill=blue!4,align=center,text width=6.5cm,inner xsep=9pt,inner ysep=7pt},
title/.style={font=\fontsize{10.5}{14}\selectfont\bfseries,text width=15cm,align=center},
label/.style={font=\fontsize{8.4}{11}\selectfont,align=center},
note/.style={label,text width=6.1cm}}
\begin{document}
\begin{tikzpicture}[x=1cm,y=1cm]
\path[use as bounding box] (-8,-3.5) rectangle(8,3.65);
\node[title] at (0,3.25) {M4-065\quad 环中的支路值与端口测量值不同};
\begin{scope}[shift={(-4.05,0)}]
\draw[wire,draw=red](1.2246467991473532e-16,2.0)--(0.6278684177437182,1.6375) (1.1041823898251595,1.3625)--(1.7320508075688774,0.9999999999999999);\begin{scope}[shift={(0.8660254037844388,1.5)},rotate=-30.000000000000004]\draw[draw=red,fill=white,line width=.8pt](-0.275,-.085)rectangle(0.275,.085);\end{scope}
\draw[wire,draw=ink](1.7320508075688774,0.9999999999999999)--(1.7320508075688774,0.275) (1.7320508075688774,-0.275)--(1.7320508075688774,-0.9999999999999999);\begin{scope}[shift={(1.7320508075688774,0.0)},rotate=-90.0]\draw[draw=ink,fill=white,line width=.8pt](-0.275,-.085)rectangle(0.275,.085);\end{scope}
\draw[wire,draw=ink](1.7320508075688774,-0.9999999999999999)--(1.1041823898251595,-1.3625) (0.6278684177437182,-1.6375)--(1.2246467991473532e-16,-2.0);\begin{scope}[shift={(0.8660254037844388,-1.5)},rotate=-150.0]\draw[draw=ink,fill=white,line width=.8pt](-0.275,-.085)rectangle(0.275,.085);\end{scope}
\draw[wire,draw=ink](1.2246467991473532e-16,-2.0)--(-0.6278684177437179,-1.6375) (-1.1041823898251593,-1.3625)--(-1.7320508075688774,-0.9999999999999999);\begin{scope}[shift={(-0.8660254037844386,-1.5)},rotate=150.0]\draw[draw=ink,fill=white,line width=.8pt](-0.275,-.085)rectangle(0.275,.085);\end{scope}
\draw[wire,draw=ink](-1.7320508075688774,-0.9999999999999999)--(-1.7320508075688772,-0.27499999999999986) (-1.7320508075688772,0.2750000000000002)--(-1.7320508075688772,1.0000000000000002);\begin{scope}[shift={(-1.7320508075688772,1.6653345369377348e-16)},rotate=90.0]\draw[draw=ink,fill=white,line width=.8pt](-0.275,-.085)rectangle(0.275,.085);\end{scope}
\draw[wire,draw=ink](-1.7320508075688772,1.0000000000000002)--(-1.104182389825159,1.3625) (-0.6278684177437178,1.6375)--(1.2246467991473532e-16,2.0);\begin{scope}[shift={(-0.8660254037844385,1.5)},rotate=29.999999999999993]\draw[draw=ink,fill=white,line width=.8pt](-0.275,-.085)rectangle(0.275,.085);\end{scope}
\fill[red] (1.2246467991473532e-16,2.0)circle(2pt);\node[above=4pt] at(1.2246467991473532e-16,2.0){$A$};
\fill[red] (1.7320508075688774,0.9999999999999999)circle(2pt);\node[above=4pt] at(1.7320508075688774,0.9999999999999999){$B$};
\fill[ink] (1.7320508075688774,-0.9999999999999999)circle(2pt);
\fill[ink] (1.2246467991473532e-16,-2.0)circle(2pt);
\fill[ink] (-1.7320508075688774,-0.9999999999999999)circle(2pt);
\fill[ink] (-1.7320508075688772,1.0000000000000002)circle(2pt);
\node[note] at (0,-2.75) {示意 $N=6$ 的一般环网络};
\end{scope}
\node[card,draw=blue!65,fill=blue!4] at (3.9,1.8) {$R_e=(N-1)r/N$};
\node[card,draw=green!65,fill=green!4] at (3.9,0) {$N\text{条边}\times R_e/r$};
\node[card,draw=orange!65,fill=orange!4] at (3.9,-1.8) {$\sum_eR_e/r=N-1$};
\end{tikzpicture}\end{document}